\begin{frame}
    \frametitle{\problemtitle}
    \begin{block}{Problem}
        You want to withdraw €$n$ ($1 \leq n \leq 10\,000$) from a cash machine. It can dispense cash in the form of $1,2,5,10,20,50,100,200,500$ (in euros).
        Check if the machine can dispense cash adding up to €$n$ which cannot be split evenly into two piles of €$n/2$.
    \end{block}
    \pause
    \begin{block}{Naive solution (Subset Sum DP + Bruteforce)}
        Given a division into coins and notes summing to €$n$, classic subset sum dynamic programming can be used to check splittability:
            $$DP[\text{considering notes up to position } i][\text{sum of cash so far}] = \text{can choose subset with this sum}$$
        Because this DP uses True / False values, it can be sped up using bitsets, to achieve a complexity of $\mathcal O( \text{\# notes} \frac{n}{w})$, where $w$ is the wordsize, typically $32$ or $64$.

        \pause
        Try out all possible splittings into all the different coins and notes, and check using the knapsack DP.

            This is too slow.
    \end{block}
\end{frame}

\begin{frame}
    \frametitle{\problemtitle}
    \begin{block}{Problem}
        % TODO: Remove this comment when you're done writing the solutions.
        You want to withdraw €$n$ ($1 \leq n \leq 10\,000$) from a cash machine. It can dispense cash in the form of $1,2,5,10,20,50,100,200,500$ (in euros).
        Check if the machine can dispense cash adding up to €$n$ which cannot be split evenly into two piles of €$n/2$.
    \end{block}
    \begin{block}{Silly solution}
        Optimize the bruteforce as much as possible, using heuristics, and precompute all the answers on your own machine.

        Then submit a big file which stores all possible answers up to $10\,000$ in it.
    \end{block}
\end{frame}

\begin{frame}
    \frametitle{\problemtitle}
    \begin{block}{Observations}
    \begin{itemize}
        \item
            For €$1$, €$5$, €$10$, €$50$ and €$100$, using $2$ or more of the same coin or note is unnecessary. \\
            For example, switching $\text{€}10 \times 2 \implies \text{€}20$ is strictly better.
        \pause
        \item
            For €$2$ and €$20$, using $5$ or more of the same coin or note is unnecessary. \\
            Switching $\text{€}2 \times 5 \implies \text{€}10$ is also strictly better.
    \end{itemize}
    \end{block}
    \pause
    \begin{block}{Easy solution}
        Run the bruteforce, but use the above observations to reduce the number of options to check.

        A rough upper bound for the number of options is $2^5 5^2 (10\,000/500) = 16\,000$, but in reality way less options actually sum to $n$.
        Here we bruteforce how many of the lower denominations we choose, also bruteforce the number of €$500$ notes.
    \end{block}
    % \pause
    % \begin{block}{Optimization}
    %     Instead of running the knapsack at the end, we can update the knapsack DP during recursive bruteforce as we add notes.
    % \end{block}
\end{frame}

\begin{frame}
    \frametitle{\problemtitle}
    \begin{block}{Greedy solution}
        It is almost possible to solve this problem greedily, but there are tricky corner cases. This is left as an exercise to the reader.
    \end{block}
    \pause
    \solvestats
\end{frame}
